\section{Криптографические основы многосторонних конфиденциальных вычислений}
\label{sec:cryptographic-foundations}

\subsection{Разделение секрета}

\subsubsection{Определение и свойства}

Разделение секрета --- это криптографическая схема, позволяющая разбить секрет \(S\) на \(n\) частей (долей) \(S_1,\ldots,S_n\) так, что любой набор не менее \(k\) долей позволяет восстановить \(S\), а из меньшего числа долей невозможно получить какую-либо информацию о секрете. Для такой \((k,n)\)-пороговой схемы характерны два свойства: во-первых, знание любых \(k\) или более долей позволяет восстановить секрет; во-вторых, знание не более \(k-1\) долей не уменьшает неопределённость секрета. Схема считается идеальной и совершенной, если доли имеют ту же длину, что и секрет, и обладают информационно-теоретической стойкостью.

\subsubsection{Схема разделения секрета Шамира}

Наиболее известной реализацией является схема Шамира (Shamir’s secret sharing).  Пусть секрет \(S\) отображается на элемент \(a_0\) конечного поля \(\mathrm{GF}(q)\).  Для порогового значения \(k\) случайно выбирается полином степени \(k-1\)

\begin{equation}
f(x) = a_0 + a_1x + a_2x^2 + \cdots + a_{k-1}x^{k-1},
\end{equation}

где \(a_1,\ldots,a_{k-1}\) --- случайные коэффициенты из \(\mathrm{GF}(q)\).  Для \(i=1,\ldots,n\) вычисляются значения \(f(i)\) и образуют пары \((i, f(i))\); каждому участнику выдаётся одна такая пара.  Любой набор из \(k\) пар позволяет восстановить \(a_0\) с помощью формулы Лагранжа:

\begin{equation}
a_0 = f(0) = \sum_{j=0}^{k-1} y_j \prod_{\substack{m=0\\m \neq j}}^{k-1} \frac{x_m}{x_m - x_j},
\end{equation}

где \((x_j, y_j)\) --- выбранные точки полинома.  Эта формула иллюстрирует, что знание \(k\) различных точек однозначно определяет полином степени \(k-1\) и, следовательно, секрет.

\subsubsection{Пример}

Рассмотрим (3,6)-пороговую схему: секрет \(S=1234\) кодируется полиномом \(f(x)=1234+166x+94x^2\), полученным выбором случайных коэффициентов 166 и 94. Из полинома вычисляются шесть долей \(f(1)=1494\), \(f(2)=1942\), \ldots, по одной для каждого из шести участников. Любые три доли позволяют восстановить секрет \(f(0)=1234\) с помощью интерполяции Лагранжа.

\subsubsection{Аддитивное и гомоморфное деление секрета}

Помимо схемы Шамира используется простое аддитивное разделение, при котором секрет разбивается на \(n\) случайных элементов поля так, что их сумма равна \(S\).  Например, для секрета \(s\) выбираются \(n-1\) случайных чисел \(r_1,\ldots,r_{n-1}\) и вычисляется \(r_n = s - \sum_{i=1}^{n-1} r_i\).  В дальнейшем операции над секретами превращаются в поэлементные операции над долями: суммируя доли по модулю \(q\), участники получают долю суммы исходных секретов.  Этот подход лежит в основе протоколов SPDZ и BMR, где для мультипликативных операций используются так называемые тройки Бивера (Beaver triples) \cite{06}.  Гомоморфное деление секрета (homomorphic secret sharing) сочетает идеи деления и шифрования и позволяет выполнять локальные операции над долями без дальнейшей коммуникации.

\subsection{Искаженные схемы (garbled circuits)}

Зашифрованная схема --- криптографический протокол, позволяющий двум сторонам вычислять любую булеву функцию от своих входных данных без раскрытия этих данных.  Протокол был предложен Яо и требует, чтобы вычисляемая функция была представлена в виде булевой схемы с двумя входами.

Пусть Алиса (garbler) и Боб (evaluator) хотят вычислить значение функции \(f(x_A, x_B)\).  Протокол состоит из следующих шагов:

\begin{itemize}
\item Построение схемы: участники согласуют общую булеву схему (набор логических элементов И, ИЛИ и пр.), реализующую функцию \(f\).

\item Гарболизация: Алиса для каждого логического провода \(w\) схемы генерирует две случайные \(k\)-битовые метки \(X^w_0, X^w_1\) для логических значений 0 и 1. Затем она проходит по всем элементам и заменяет значения в таблице истинности соответствующими метками. После этого она шифрует каждую строку таблицы двойным симметричным шифрованием с ключами, соответствующими двум входным меткам, и случайно переставляет строки; результатом является искаженная таблица.

\item Передача схемы и меток: Алиса отправляет Бобу зашифрованные таблицы всех элементов и метки, соответствующие своим входам. Метки Боба для каждого бита он получает через протокол слепой передачи (см. п. 2.4), что гарантирует, что Алиса не узнает его выбор.

\item Оценивание: Боб, имея искаженную схему и свои метки, проходит по элементам схемы, подбирая для каждого элемента единственную строку таблицы, которую он может расшифровать. В результате он получает метку на выходном логическом проводе и декодирует её с помощью соответствующего алгоритма, получая значение функции.
\end{itemize}

В формальном описании схема искажения задается рандомизированным алгоритмом \(\text{Gb}\), который для функции \(f\) формирует тройку \((F,e,d)\): \(F\) --- искаженное представление функции, \(e\) --- функция кодирования входов, а \(d\) --- функция декодирования выхода. При вычислении \(F(e(x))\) получается искаженный выход, который после применения \(d\) преобразуется в значение \(f(x)\).

\subsection{Гомоморфное шифрование (homomorphic encryption)}

\subsubsection{Определение и свойства}

Гомоморфное шифрование --- это криптографический метод, позволяющий выполнять вычислительные операции непосредственно над зашифрованными данными без их предварительного расшифрования. При этом расшифрование результата вычислений позволяет получить значение, эквивалентное результату выполнения соответствующих операций над исходными данными. Возможность выполнения вычислений заданной глубины определяется особенностями криптографической схемы и выбранными параметрами. Для обеспечения вычислений большей глубины в ряде схем применяется процедура бутстрэппинга, позволяющая уменьшать накопленный в процессе вычислений шум в шифротекстах.

\subsubsection{Частично гомоморфное шифрование (PHE)}

В простейших схемах выполняется только одно арифметическое действие.  Например, в криптосистеме RSA (без паддинга) шифротекст \(E(m) = m^e \bmod n\) обладает мультипликативным свойством: произведение двух шифротекстов равно шифротексту произведения исходных сообщений.  Схема Эль-Гамаля использует пары \((g^r, m \cdot h^r)\), и перемножение двух шифротекстов приводит к шифрованию произведения сообщений.  В схеме Паилье (Paillier) сообщение шифруется как \(E(m) = g^m r^n \bmod n^2\), и произведение шифротекстов даёт шифр суммы: \(E(m_1)\cdot E(m_2) = E(m_1 + m_2)\).  Эти свойства позволяют строить МКВ-протоколы для простых операций без интерактивности.

\subsubsection{Полностью гомоморфное шифрование (FHE)}

Полностью гомоморфные схемы выполняют как сложение, так и умножение над шифротекстами.  Идея была предложена в 1978 году, но первый практически реализуемый FHE-протокол разработал Гентри в 2009 году, введя метод «bootstrap» для подавления шума.  Схема строится на «частично» гомоморфном шифре, который поддерживает ограниченное число операций, и периодическом гомоморфном декодировании, уменьшающем шум.  Теоретические описания FHE используют абстракцию колец и операции AND/XOR как аналоги умножения/сложения.  Существующие реализации (CKKS, BGV, BFV) основаны на сложности задач «обучения с ошибками» и «приближенного НОД»; они позволяют вычислять целые и вещественные функции над зашифрованными данными, но требуют значительных вычислительных ресурсов и больших размеров ключей.

\subsection{Протокол слепой передачи (oblivious transfer)}

Определение.  Ослеплённая передача (oblivious transfer, OT) --- это криптографический протокол, в котором отправитель имеет набор сообщений, а получатель получает ровно одно из них, причём отправитель остаётся в неведении о выборе получателя.  Первая версия протокола была предложена М. Рабином, где отправитель посылал сообщение с вероятностью 1/2, и не знал, получил ли его адресат.  Более полезная 1-из-2 OT-схема, разработанная Эвеном, Голдрайхом и Лемпелем, позволяет получателю выбирать один из двух сообщений, не раскрывая свой выбор, и обобщается до 1-из-n OT, где клиент получает один элемент базы данных, ничего не узнавая о других элементах.  Ослеплённая передача является полной для безопасных вычислений: имея реализацию OT, можно безопасно вычислить любую полиномиально вычислимую функцию.

Пример протокола 1-из-2 на основе RSA:
\begin{equation}
\begin{aligned}
v &= (x_b + k^e)\bmod N, \\
m'_i &= \bigl(m_i + (v-x_i)^d\bigr)\bmod N, \quad i\in\{0,1\}, \\
m_b &= (m'_b-k)\bmod N.
\end{aligned}
\end{equation}

Здесь \(N,e,d\) --- параметры RSA, \(x_0,x_1,k\) --- случайные значения, а \(b\in\{0,1\}\) --- выбор Боба, который восстанавливает \(m_b\), не раскрывая Алисе свой выбор и не получая сведений о \(m_{1-b}\).

Роль в МКВ. OT применяется для безопасной передачи входных меток от garbler’а к evaluator’у и в функциональных протоколах, включая Private Set Intersection и oblivious PRF. Возможность реализации произвольных МКВ-протоколов следует из полноты OT.

\subsection{Сравнительный анализ криптографических примитивов}

В таблице~\ref{tab:mpc-primitives} сопоставлены криптографические примитивы МКВ; выбор зависит от модели угроз, числа участников, вычисляемой функции и допустимых затрат.

{\small
\setlength{\tabcolsep}{3pt}
\renewcommand{\arraystretch}{1.2}
\setlength{\tablecontentwidth}{\dimexpr\textwidth-10\tabcolsep-6\arrayrulewidth\relax}
\begin{longtable}{|P{0.16\tablecontentwidth}|P{0.20\tablecontentwidth}|P{0.22\tablecontentwidth}|P{0.22\tablecontentwidth}|P{0.20\tablecontentwidth}|}
\caption{Сравнение криптографических примитивов МКВ}\label{tab:mpc-primitives}\\
\hline
\multicolumn{1}{|H{0.16\tablecontentwidth}|}{\textbf{Примитив}} & \multicolumn{1}{H{0.20\tablecontentwidth}|}{\textbf{Основная идея}} & \multicolumn{1}{H{0.22\tablecontentwidth}|}{\textbf{Преимущество}} & \multicolumn{1}{H{0.22\tablecontentwidth}|}{\textbf{Недостаток}} & \multicolumn{1}{H{0.20\tablecontentwidth}|}{\textbf{Область применения}} \\
\hline
\endfirsthead
\multicolumn{5}{@{}l@{}}{\normalsize Продолжение таблицы~\thetable}\\
\hline
\multicolumn{1}{|H{0.16\tablecontentwidth}|}{\textbf{Примитив}} & \multicolumn{1}{H{0.20\tablecontentwidth}|}{\textbf{Основная идея}} & \multicolumn{1}{H{0.22\tablecontentwidth}|}{\textbf{Преимущество}} & \multicolumn{1}{H{0.22\tablecontentwidth}|}{\textbf{Недостаток}} & \multicolumn{1}{H{0.20\tablecontentwidth}|}{\textbf{Область применения}} \\
\hline
\endhead
Разделение секрета & Секрет кодируется долями; \(k\) из \(n\) долей позволяют восстановить секрет. & Информационно-теоретическая стойкость; лёгкие арифметические операции; эффективная реализация для линейных функций. & Требует много раундов обмена сообщениями для умножения; каждый участник должен быть онлайн. & Протоколы SPDZ/BMR, защищённое машинное обучение, приватная аналитика данных. \\
\hline
Garbled circuits & Функция представляется булевой схемой, её таблицы истинности шифруются случайными метками. & Небольшое число раундов; поддержка произвольных функций; удобство для сравнений. & Большие таблицы и объём передачи; обычно применяется для двух участников; криптографические операции на каждом элементе. & Задача миллионеров Яо, приватное сравнение, доказательства с нулевым разглашением. \\
\hline
Гомо\-морфное шифрование & Операции выполняются над шифротекстами; после расшифрования результат совпадает с результатом операции над открытыми данными. & Отсутствие интерактивности; одна сторона может выполнить вычисления, не раскрывая входы. & Высокие вычислительные затраты; ограниченная глубина при FHE \cite{02}; большие ключи. & Облачные вычисления, приватные ML-сервисы, вычисления на единичных данных. \\
\hline
Oblivious transfer & Получатель выбирает и получает одно из \(n\) сообщений, не раскрывая выбора; отправитель не узнаёт, какое сообщение выбрано. & Универсальность и полнота для МКВ; удобство для ввода скрытых данных. & Сам по себе не вычисляет функцию; требует криптографических операций, например RSA или протокола Диффи--Хеллмана. & Ввод входов в garbled circuits, PSI (Private Set Intersection), oblivious PRF. \\
\hline
\end{longtable}
}

\subsection{Вывод}

Пороговые схемы разделения секрета обеспечивают информационно-теоретическую защиту и восстановление секрета по долям (\(k\)-кворум); характеристики остальных примитивов сопоставлены в таблице~\ref{tab:mpc-primitives}. В следующем разделе протоколы классифицированы по модели противника, числу участников и криптографическим механизмам.